% GENERATED FILE. Edit canonical lessons, then run npm run generate:books.
% canonical-source-hash: fnv1a64:608e13128b9b58f4
% canonical-lessons: GU-C550-tir, GU-C550-nisan, GU-C550-taakvu, GU-C550-bhedvu, GU-C550-dodvir

\chapter{Taking Aim}
\label{ch:gu-taking-aim}
\hlchaptermodality{550}

\begin{chapteropening}
\textbf{By the end of this chapter:} \emph{I can talk about aiming at a target and hitting it.}

Complete the last lesson of chapter 550: I can talk about aiming at a target and hitting it.
\end{chapteropening}

\section[\texorpdfstring{\gu{તીર} (tīr)}{તીર (tīr)}]{\gu{તીર} (tīr) — an arrow}

\label{lesson:GU-C550-tir}



\begin{quote}
Before the new one: say the Gujarati for to catch a ball, then the Gujarati for a medal.
\end{quote}

\subsection*{You'll want to know: \gu{તીર}}
\textbf{\gu{તીર}} — \emph{tīr} — \textquotedblleft{}an arrow\textquotedblright{}.

\begin{quote}
\gu{જો} \gu{પવન} \gu{નહીં} \gu{હોય} \gu{તો} \gu{તીર} \gu{સીધું} \gu{જશે}.
\end{quote}

\begin{quote}
\emph{jo pavan nahī̃ hoy to tīr sīdhũ jaśe.}
\end{quote}

\begin{quote}
If there's no wind, the arrow will fly straight.
\end{quote}

A neuter noun.

\begin{cousinweb}
\textbf{\gu{તીર}} (\emph{tīr}) is usually traced to Persian \emph{tīr}, \textquotedblleft{}an arrow\textquotedblright{}.
\end{cousinweb}

\begin{grammarlens}[title={What you've built}]
One more word for saying what happens if.
\end{grammarlens}

\subsection*{Your turn}
\begin{itemize}
\raggedright
  \item \emph{Say it:} \emph{tīr}
  \item \emph{Say it:} \emph{tīr}, once more
  \item \emph{Say it:} \emph{jo pavan nahī̃ hoy to tīr sīdhũ jaśe.} — If there's no wind, the arrow will fly straight.
  \item \emph{Recall:} say the Gujarati for to catch a ball, then the Gujarati for a medal, then say \emph{tīr} again
\end{itemize}

\begin{tcolorbox}[breakable,colback=teal!4,colframe=teal!35!black,title={Before you move on}]
What does \gu{તીર} mean? (\textquotedblleft{}An arrow\textquotedblright{}.) Say it once more.
\end{tcolorbox}
\section[\texorpdfstring{\gu{નિશાન} (niśān)}{નિશાન (niśān)}]{\gu{નિશાન} (niśān) — a target}

\label{lesson:GU-C550-nisan}



\begin{quote}
Before the new one: say the Gujarati for a medal, then the Gujarati for an arrow.
\end{quote}

\subsection*{You'll want to know: \gu{નિશાન}}
\textbf{\gu{નિશાન}} — \emph{niśān} — \textquotedblleft{}a target\textquotedblright{}.

\begin{quote}
\gu{જો} \gu{તમારું} \gu{તીર} \gu{નિશાન} \gu{પર} \gu{લાગશે} \gu{તો} \gu{તમે} \gu{જીતશો}.
\end{quote}

\begin{quote}
\emph{jo tamārũ tīr niśān par lāgśe to tame jītśo.}
\end{quote}

\begin{quote}
If your arrow hits the target, you'll win.
\end{quote}

A neuter noun.

\begin{cousinweb}
\textbf{\gu{નિશાન}} (\emph{niśān}) is usually traced to Persian \emph{nishān}, \textquotedblleft{}a mark, a sign\textquotedblright{}; the mark one aims at is a target.
\end{cousinweb}

\begin{grammarlens}[title={What you've built}]
One more word for saying what happens if.
\end{grammarlens}

\subsection*{Your turn}
\begin{itemize}
\raggedright
  \item \emph{Say it:} \emph{niśān}
  \item \emph{Say it:} \emph{niśān}, once more
  \item \emph{Say it:} \emph{jo tamārũ tīr niśān par lāgśe to tame jītśo.} — If your arrow hits the target, you'll win.
  \item \emph{Recall:} say the Gujarati for a medal, then the Gujarati for an arrow, then say \emph{niśān} again
\end{itemize}

\begin{tcolorbox}[breakable,colback=teal!4,colframe=teal!35!black,title={Before you move on}]
What does \gu{નિશાન} mean? (\textquotedblleft{}A target\textquotedblright{}.) Say it once more.
\end{tcolorbox}
\section[\texorpdfstring{\gu{તાકવું} (tākvũ)}{તાકવું (tākvũ)}]{\gu{તાકવું} (tākvũ) — to aim at}

\label{lesson:GU-C550-taakvu}



\begin{quote}
Before the new one: say the Gujarati for an arrow, then the Gujarati for a target.
\end{quote}

\subsection*{You'll want to know: \gu{તાકવું}}
\textbf{\gu{તાકવું}} — \emph{tākvũ} — \textquotedblleft{}to aim at\textquotedblright{}.

\begin{quote}
\gu{જો} \gu{તમે} \gu{ધ્યાનથી} \gu{તાકશો} \gu{તો} \gu{તીર} \gu{નિશાન} \gu{પર} \gu{લાગશે}.
\end{quote}

\begin{quote}
\emph{jo tame dhyānthī tākśo to tīr niśān par lāgśe.}
\end{quote}

\begin{quote}
If you aim carefully, the arrow will hit the target.
\end{quote}

Past, with \gu{તેણે} (he or she …): \gu{તેણે} \gu{નિશાન} \gu{તાક્યું}, \emph{teṇe niśān tākyũ}.

\begin{cousinweb}
\textbf{\gu{તાકવું}} (\emph{tākvũ}) is filed by CDIAL under Sanskrit \emph{tarkáyati}, \textquotedblleft{}considers\textquotedblright{}, and CDIAL lists the Gujarati verb, \textquotedblleft{}to look at\textquotedblright{} (CDIAL 5716); one looks hard at what one aims at.
\end{cousinweb}

\begin{grammarlens}[title={What you've built}]
One more word for saying what happens if.
\end{grammarlens}

\subsection*{Your turn}
\begin{itemize}
\raggedright
  \item \emph{Say it:} \emph{tākvũ}
  \item \emph{Say it:} \emph{tākvũ}, once more
  \item \emph{Say it:} \emph{jo tame dhyānthī tākśo to tīr niśān par lāgśe.} — If you aim carefully, the arrow will hit the target.
  \item \emph{Recall:} say the Gujarati for an arrow, then the Gujarati for a target, then say \emph{tākvũ} again
\end{itemize}

\begin{tcolorbox}[breakable,colback=teal!4,colframe=teal!35!black,title={Before you move on}]
What does \gu{તાકવું} mean? (\textquotedblleft{}To aim at\textquotedblright{}.) Say it once more.
\end{tcolorbox}
\section[\texorpdfstring{\gu{ભેદવું} (bhedvũ)}{ભેદવું (bhedvũ)}]{\gu{ભેદવું} (bhedvũ) — to pierce (a target)}

\label{lesson:GU-C550-bhedvu}



\begin{quote}
Before the new one: say the Gujarati for a target, then the Gujarati for to aim at.
\end{quote}

\subsection*{You'll want to know: \gu{ભેદવું}}
\textbf{\gu{ભેદવું}} — \emph{bhedvũ} — \textquotedblleft{}to pierce (a target)\textquotedblright{}.

\begin{quote}
\gu{જો} \gu{અરુણ} \gu{પહેલા} \gu{તીરથી} \gu{નિશાન} \gu{ભેદત} \gu{તો} \gu{તે} \gu{જીતી} \gu{જાત}.
\end{quote}

\begin{quote}
\emph{jo Aruṇ pahelā tīrthī niśān bhedat to te jītī jāt.}
\end{quote}

\begin{quote}
If Arun had hit the target with the first arrow, he would have won.
\end{quote}

Past, with \gu{તેણે} (he or she …): \gu{તેણે} \gu{નિશાન} \gu{ભેદ્યું}, \emph{teṇe niśān bhedyũ}.

\begin{cousinweb}
\textbf{\gu{ભેદવું}} (\emph{bhedvũ}) goes back to Sanskrit \emph{bhid}, \textquotedblleft{}to split, to pierce\textquotedblright{} (Monier-Williams).
\end{cousinweb}

\begin{grammarlens}[title={What you've built}]
One more word for saying what happens if.
\end{grammarlens}

\subsection*{Your turn}
\begin{itemize}
\raggedright
  \item \emph{Say it:} \emph{bhedvũ}
  \item \emph{Say it:} \emph{bhedvũ}, once more
  \item \emph{Say it:} \emph{jo Aruṇ pahelā tīrthī niśān bhedat to te jītī jāt.} — If Arun had hit the target with the first arrow, he would have won.
  \item \emph{Recall:} say the Gujarati for a target, then the Gujarati for to aim at, then say \emph{bhedvũ} again
\end{itemize}

\begin{tcolorbox}[breakable,colback=teal!4,colframe=teal!35!black,title={Before you move on}]
What does \gu{ભેદવું} mean? (\textquotedblleft{}To pierce (a target)\textquotedblright{}.) Say it once more.
\end{tcolorbox}
\section[\texorpdfstring{\gu{દોડવીર} (doḍvīr)}{દોડવીર (doḍvīr)}]{\gu{દોડવીર} (doḍvīr) — a runner (an athlete)}

\label{lesson:GU-C550-dodvir}



\begin{quote}
Before the new one: say the Gujarati for to aim at, then the Gujarati for to pierce (a target).
\end{quote}

\subsection*{You'll want to know: \gu{દોડવીર}}
\textbf{\gu{દોડવીર}} — \emph{doḍvīr} — \textquotedblleft{}a runner (an athlete)\textquotedblright{}.

\begin{quote}
\gu{જો} \gu{આ} \gu{દોડવીર} \gu{રોજ} \gu{દોડશે} \gu{તો} \gu{તે} \gu{ચંદ્રક} \gu{જીતશે}.
\end{quote}

\begin{quote}
\emph{jo ā doḍvīr roj doḍśe to te candrak jītśe.}
\end{quote}

\begin{quote}
If this runner runs every day, she will win a medal.
\end{quote}

A masculine or feminine noun, by the person meant.

\begin{cousinweb}
\textbf{\gu{દોડવીર}} (\emph{doḍvīr}) is built from the stem of \textbf{\gu{દોડવું}}, \textquotedblleft{}to run\textquotedblright{}, which CDIAL files under Sanskrit \emph{drávati} (CDIAL 6624), and \emph{vīr}, \textquotedblleft{}a hero\textquotedblright{}, from Sanskrit \emph{vīra} (Monier-Williams).
\end{cousinweb}

\begin{grammarlens}[title={What you've built}]
One more word for saying what happens if.
\end{grammarlens}

\subsection*{Your turn}
\begin{itemize}
\raggedright
  \item \emph{Say it:} \emph{doḍvīr}
  \item \emph{Say it:} \emph{doḍvīr}, once more
  \item \emph{Say it:} \emph{jo ā doḍvīr roj doḍśe to te candrak jītśe.} — If this runner runs every day, she will win a medal.
  \item \emph{Recall:} say the Gujarati for to aim at, then the Gujarati for to pierce (a target), then say \emph{doḍvīr} again
\end{itemize}

\begin{tcolorbox}[breakable,colback=teal!4,colframe=teal!35!black,title={Before you move on}]
What does \gu{દોડવીર} mean? (\textquotedblleft{}A runner (an athlete)\textquotedblright{}.) Say it once more.
\end{tcolorbox}
